\begin{proposition}[The lower endpoint for abelian-type coefficients]
\label{c:abelian-invariant-weights}
\status{Proved}
Let $K$ be a complete discretely valued field with finite residue field of
characteristic $p$, and let $\ell\ne p$.  Suppose that $V$ is the
$\ell$-adic realization of a pure weight-$j$ motive which, after a finite
extension of $K$, is a direct summand of tensor constructions on the
$H^1$ of abelian varieties, their duals and Tate twists.  Finite-image
Artin factors are allowed.  Then the eigenvalues of geometric Frobenius
on $V^{I_K}$ have weights at most $j$.  In particular this holds for the
local realizations of pure absolute-Hodge or motivated motives of abelian
type over number fields.  The bound is attained for good reduction.
\end{proposition}

\begin{proof}
We give the weight--monodromy argument, including the tensor step.
After a finite extension we may suppose that all the finitely many
abelian generators have semistable reduction and all the Artin factors
are trivial.  Semistable reduction is SGA~7~I, IX, Theorem~3.6, printed
p.~351.  For one such abelian variety write
$U=T_\ell(A)\otimes\Q_\ell$, $V_0=U^{I_K}$, and let $W$ be the
toric part of its Tate module.  The N\'eron special fibre gives
\[
 0\longrightarrow W\longrightarrow V_0
       \longrightarrow T_\ell(B)\otimes\Q_\ell\longrightarrow0,
\]
where $B$ is the abelian quotient of the identity component.  The exact
sequence and the identification of the fixed part are IX, formula
(2.1.8), p.~328, and Proposition~2.2.5, pp.~332--333.  A torus has Tate
module of weight $-2$, and $T_\ell(B)$ has weight $-1$.  A polarization
and the orthogonality theorem IX, Theorem~2.4, p.~338, give
$W=V_0^\perp$ in the semistable case, as stated in Proposition~3.5,
p.~350.  Consequently $U/V_0$ is dual to $W$ with values in
$\Q_\ell(1)$ and has weight $0$.

Corollary~3.5.2, p.~351, says that inertia factors through the tame
$\Z_\ell(1)$ quotient and that $(g-1)^2=0$ for a generator $g$.
Thus its logarithm $N$ satisfies $N^2=0$ and $\ker N=V_0$.
The polarization is inertia invariant, so $N$ is skew adjoint and
$\operatorname{im}N=(\ker N)^\perp=W$.  It induces an isomorphism
\[
 N:U/V_0\xrightarrow{\sim}W(-1).
\]
All these statements are rational statements; no assumption on an
integral monodromy pairing is required.  The filtration
$0\subset W\subset V_0\subset U$ is therefore the monodromy
filtration centred at weight $-1$, with graded weights $-2,-1,0$.
Dualizing proves weight--monodromy, centred at $1$, for $H^1(A)$.
This derives exactly the form we use from the cited original pages
of SGA~7 \citep{c_sga7}.

For completeness, pure Weil--Deligne representations are stable under
the indicated tensor operations.  The monodromy filtration is
characterized by $NM_a\subset M_{a-2}(-1)$ and the isomorphisms
$N^a:\operatorname{gr}_a^M\xrightarrow{\sim}
\operatorname{gr}_{-a}^M(-a)$ for $a\geq0$.  On a tensor product
the convolution filtration has these properties: decompose the
nilpotent modules into $\mathfrak{sl}_2$ strings and apply the
Clebsch--Gordan decomposition to their tensor products.  Its graded
terms have the sums of the two graded weights.  Thus the centre
weights add.  Duality reverses the indices and weights, and a Tate
twist shifts the centre weight by an even integer.  An idempotent
commuting with the Weil action and $N$ preserves the unique
monodromy filtration; its image on every graded term is pure of the
same weight.  This proves stability under direct summands as well.
We apply these facts to the weight-$j$ part of the tensor presentation
of the given pure motive.  Its Weil--Deligne representation is pure
of weight $j$.

For a monodromy string centred at $j$, the vector killed by $N$ is
its lowest-weight vector.  Equivalently, $N$ is injective on every
positive-index graded term of its monodromy filtration.  Hence
$\ker N$ has weights at most $j$.  The original inertia invariants
are a Frobenius-stable subspace of $\ker N$ after the finite
extension.  If that extension has residue degree $f$, any lift of
its geometric Frobenius differs from the $f$th power of the original
Frobenius by an element of $I_K$.  These two operators therefore
agree on $V^{I_K}$.  The weight inequality is unchanged
on taking these powers, and therefore descends to $K$.
For a pure motive with good reduction the invariant space is the
whole pure weight-$j$ representation, proving the stated sharpness.
\end{proof}

